\begin{pastexam}[label={px:gk3-111-ds-02}]{高考三級資訊處理 111 年 第二題}
（申論；25 分）

利用鏈結串列（Linked list）實做佇列（Queues），給予如下鏈結串列節點及佇列定義，\texttt{front}
指標指在串列第一個節點，rear 指標指在串列最後一個節點，請使用 C 語言完成 \texttt{insert(pq, x)}程序，
將整數值 x 加入（\texttt{Insert}）到佇列，程式需檢查佇列加入前是否為空的鏈結串列，可使用函數
getnode() 配置（Allocate）一新節點。（25 分）
\begin{minted}[fontsize=\small, frame=single]{c}
struct node{
    int info;
    struct node *next;
};
typedef struct node *NODEPTR;

struct queue{
    NODEPTR  front, rear;
};
struct queue q;

NODEPTR  getnode()
{
   NODEPTR  p;
   p = (NODEPTR)malloc(sizeof(struct node));
   return(p);
}

insert(pq, x)
struct queue *pq;
int x;
{
     NODEPTR p;

}
\end{minted}
\end{pastexam}

\begin{pastsoln}
\textbf{想法。} 新節點一定接在串列尾端（佇列從尾端加入、從前端取出）。 分兩種情況：
\begin{itemize}
\item 加入前佇列是空的（\texttt{pq->rear == NULL}）：新節點是唯一的節點，\texttt{front} 和
\texttt{rear} 都要指向它。
\item 不是空的：把新節點接在目前最後一個節點後面（\texttt{pq->rear->next = p}），再讓
\texttt{rear} 指向它；\texttt{front} 不變。
\end{itemize}
兩種情況最後都要做 \texttt{pq->rear = p}，新節點的 \texttt{next} 都設成 \texttt{NULL}。
這就是資結書虛擬碼~6.5 的串列佇列加入。 題目給的是舊式（K\&R）的參數宣告寫法，照原樣補完：

\begin{code}{c}{insert：把 x 加到佇列尾端}
insert(pq, x)
struct queue *pq;
int x;
{
     NODEPTR p;

     p = getnode();
     if (p == NULL) {                 /* 配置失敗 */
         printf("queue overflow\n");
         exit(1);
     }
     p->info = x;
     p->next = NULL;                  /* 新節點是最後一個 */
     if (pq->rear == NULL)            /* 加入前是空佇列 */
         pq->front = p;
     else
         pq->rear->next = p;          /* 接在原本最後一個後面 */
     pq->rear = p;
}
\end{code}

\textbf{說明。}
\begin{itemize}
\item 空佇列的判斷用 \texttt{rear == NULL}（也可以用 \texttt{front == NULL}），前提是佇列初始化
和刪除最後一個節點時，\texttt{front}、\texttt{rear} 都設成 \texttt{NULL}。
\item 順序很重要：要先把 \texttt{pq->rear->next} 指向新節點，再移動 \texttt{rear}。 先移動
\texttt{rear} 就找不到原本的最後一個節點了。
\item 空佇列時如果漏掉 \texttt{front = p}，之後取出會從 \texttt{NULL} 開始，佇列永遠是空的；
如果沒判斷就直接寫 \texttt{pq->rear->next}，會對 \texttt{NULL} 取值。 這就是題目要求檢查的原因。
\item \texttt{getnode()} 內用 \texttt{malloc}，可能回傳 \texttt{NULL}，所以先檢查。
\end{itemize}

\textbf{追蹤。} 空佇列依序加入 5、8：加入 5 時 \texttt{rear} 為 \texttt{NULL}，\texttt{front} 和
\texttt{rear} 都指向 5；加入 8 時 5 的 \texttt{next} 指向 8，\texttt{rear} 改指 8，串列是
$\texttt{front} \to 5 \to 8 \leftarrow \texttt{rear}$。

\textbf{複雜度。} 只做固定幾個指標操作，不需要走過串列，時間 $\Theta(1)$。
\end{pastsoln}
